\chapter{Future Work}\label{chap:todo}

The delivered system meets all stated objectives and is fully functional as a
cross-platform desktop application.
Nevertheless, a number of directions for improvement and extension were identified
during development.
These are grouped into short-term improvements that address known gaps in the current
implementation and longer-term enhancements that would expand the scope of the system.

% ─────────────────────────────────────────────────────────────────────────────
\section{Short-Term Improvements}\label{sec:fw_shortterm}
% ─────────────────────────────────────────────────────────────────────────────

\paragraph{Sub-task backend implementation.}
The Table View renders a placeholder sub-task row for each task, but no corresponding
backend exists.
Implementing sub-tasks would require adding a \texttt{SubTask} entity with a foreign
key to \texttt{TaskItem}, a database migration, a \texttt{SubTasksController}, and a
\texttt{SubTaskRepository}.
The \texttt{ISyncableEntity} interface would then be applied to \texttt{SubTask} so
that sub-tasks participate in the offline synchronisation pipeline automatically.
The UI work --- expanding the Table View row and exposing CRUD controls --- is largely
already scaffolded.
This is the highest-priority item because it is the largest gap between the stated
scope and the delivered system.

\paragraph{Mistral API key security.}
As noted in Section~\ref{sec:concl_limitations}, the Mistral API key is currently a
compile-time constant.
The immediate remedy is to read the key from an environment variable
(\texttt{TASKFLOW\_MISTRAL\_KEY}) at startup, falling back to the
\texttt{appsettings.json} configuration value, and ultimately defaulting to the
embedded constant only as a last resort.
In a production distribution the key should be provisioned by the installer or entered
by the user on first launch and stored in the operating system's credential store
(\texttt{Windows Credential Manager}, macOS \texttt{Keychain}, or \texttt{libsecret}
on Linux).

\paragraph{Automated test suite.}
The absence of automated tests is the most significant quality risk in the codebase.
A practical test suite would consist of three layers:
\begin{enumerate}
  \item \textbf{Unit tests} (xUnit) for pure service methods and the sync engine logic,
        using mocked repository interfaces.
  \item \textbf{API integration tests} using \texttt{WebApplicationFactory<Program>} to
        spin up an in-memory ASP.NET Core server backed by an in-memory SQLite database,
        covering all controller endpoints.
  \item \textbf{End-to-end tests} using Playwright to exercise the packaged Electron
        application, verifying that key user journeys (login, create task, go offline,
        create task, go online, confirm sync) complete correctly.
\end{enumerate}
Adding this suite would allow regressions to be detected automatically and would
provide the confidence required to accelerate future feature development.

\paragraph{Per-field conflict resolution.}
The current last-write-wins strategy resolves sync conflicts at the document level:
the document with the later \texttt{UpdatedAt} timestamp overwrites the earlier one
entirely.
A more granular approach would track a per-field version vector, merging non-conflicting
field changes and presenting a conflict dialogue only when the same field was edited
concurrently.
This would be a meaningful improvement for multi-device users who edit different
properties of the same task on different machines while offline.

\paragraph{macOS notarisation.}
Submitting the macOS distribution to Apple's notarisation service would eliminate the
Gatekeeper manual-override step described in Section~\ref{sec:res_distribution}.
Notarisation requires an Apple Developer Programme membership and a CI pipeline step
that submits the signed DMG to Apple and waits for approval, but the result is a
significantly smoother first-launch experience for macOS users.

% ─────────────────────────────────────────────────────────────────────────────
\section{Long-Term Enhancements}\label{sec:fw_longterm}
% ─────────────────────────────────────────────────────────────────────────────

\paragraph{Mobile companion application.}
A React Native application sharing the MongoDB relay layer would allow users to view
and update tasks on iOS and Android devices.
Because the MongoDB layer already stores a current snapshot of every syncable entity,
the mobile client would not require an additional backend service: it could connect
directly to MongoDB (via Atlas Data API or a lightweight proxy) and write changes that
the desktop client would pick up on its next sync cycle.

\paragraph{AI-powered task creation from natural language.}
The existing Mistral integration could be extended to accept a free-text description
such as \emph{``Set up a team meeting next Thursday at 10 AM to review the sprint''}
and return a structured \texttt{CreateTaskDto} or \texttt{CreateCalendarEventDto}
payload that is immediately submitted to the appropriate API endpoint.
This would reduce the number of form fields a user must complete for routine task
creation and leverage the AI assistant for a workflow-acceleration purpose beyond
question answering.

\paragraph{CRDT-based synchronisation.}
Replacing last-write-wins with a Conflict-free Replicated Data Type (CRDT) would allow
concurrent edits from multiple devices to be merged deterministically without data
loss.
A suitable approach for the task entity would be to model each editable field as a
Last-Write-Wins Register (LWW-Register) with a logical clock, producing merge results
that are commutative, associative, and idempotent.
This would require replacing the current \texttt{UpdatedAt}-based merge in
\texttt{UserDataSyncService} with a per-field CRDT merge function, and storing the
logical clock alongside each field in the MongoDB document.

\paragraph{Analytics and reporting dashboard.}
The current dashboard provides aggregate counts and a seven-day task trend.
A richer analytics view would include burndown charts, task completion velocity per
team member, average task age by priority, and project-level progress bars.
These metrics could be computed with MongoDB aggregation pipelines over the relay
collection, avoiding additional load on the local SQLite database.

\paragraph{Third-party integration plugins.}
A plugin API that allows third-party integrations to be registered as background
services would extend the application's reach.
High-value initial integrations would include bidirectional synchronisation with
GitHub Issues (mapping issues to tasks), Google Calendar (mapping CalendarEvents to
Google Calendar entries), and Slack (posting task assignment and due-date notifications
to a channel).
Each plugin would implement a common \texttt{IIntegrationPlugin} interface, keeping
the core application decoupled from external service details.

\paragraph{Cloud-hosted multi-tenant mode.}
The current architecture is single-user and desktop-only.
A future SaaS variant would replace the embedded Electron shell with a web-hosted
frontend, the per-user SQLite database with a per-organisation PostgreSQL instance
managed by Entity Framework Core, and the MongoDB relay with a shared cluster.
The ASP.NET Core service layer would require minimal changes because its business logic
is already decoupled from the storage layer via repository interfaces;
the main effort would be in the authentication layer (multi-tenant JWT issuance,
organisation scoping on all queries) and the deployment pipeline (containerisation,
orchestration, secrets management).
